\documentclass[11pt,a4paper]{article}
\usepackage[margin=23mm]{geometry}
\usepackage{amsmath,amssymb,amsthm}
\setlength{\emergencystretch}{2em}
\newtheorem{lemma}{Lemma}
\newtheorem{theorem}[lemma]{Theorem}
\newcommand{\gen}[1]{\langle#1\rangle}
\newcommand{\Sym}{\operatorname{Sym}}
\title{A separator for unequal witness intervals}
\author{Human Compression project: L044 continuation}
\date{}
\begin{document}\maketitle
All marks below are involutions satisfying string
commutations. The restricted actions used in the
first lemma need no intersection property.

\begin{lemma}[Equality in the truncated count]
Retain the actual local crossing setting of the
truncated-orbit theorem, with active labels
$m,\ldots,y$, crossing colour $j$, and
$j+1\le x<y$. Let $L=Q_{<j}$,
$R'=Q_{\{j+1,\ldots,x-1\}}$ and
$C_A=a^{LR'}$, $C_B=b^{LR'}$ for the crossing
endpoints. Suppose
\[
 h=j-m\ge2,\quad v=x-j-1\ge1,\quad
                  |C_A|+|C_B|=h+v+2.
\]
If the first omission $Q_{>m}$ is intransitive,
it has exactly two orbits: a singleton $\{z\}$
contained in $C_A\cup C_B$, and its transitive
complement in the full local orbit $W$.
\end{lemma}
\begin{proof}
The left and truncated right active lengths have
maxima $h,v$, exactly as proved in the complete
truncated count. If $h+v\ge4$, the sharper two-factor bound
makes simultaneous attainment on the same side
impossible: that side alone would require
$h+v+2$ points. The only remaining pair is
$(h,v)=(2,1)$. If both maxima occurred on one
side, equality in the earlier bound would give
$|C_A|=4$, $|C_B|=1$, after renaming. The
single-involution factor $R'$ cannot be transitive
on four points. If both factors were intransitive,
the orbit bound would give at least six points.
Thus $L$ is transitive on $C_A$. Every later
right mark fixes $a,b$ and centralizes $L$,
so fixes $C_A$ and $C_B$ pointwise. Both sets
would therefore be invariant under the full
$LR$, making them the complete omission sides.
No mark after $j+1$ could be active, contrary
to $y>x=j+2$. This excludes simultaneous
attainment in the last case too. Rename the sides so that
\[
 |C_A|=h+1,\qquad |C_B|=v+1,
 \qquad L|_{C_A}=\Sym(C_A),\quad R'|_{C_B}=\Sym(C_B).
\]
The left degree is at least three, so commutation
makes $R'$ fix $C_A$ pointwise. Every further
right mark fixes $a$, being at distance at least
two from $j$. Thus the full right factor
$R=Q_{>j}$ fixes $a$. Since $LR$ is transitive
on the actual side $A$, we have $A=a^L=C_A$,
with $R$ fixing all of $A$.
If $v\ge2$, the right degree is also at least
three, and $L$ fixes $C_B$ pointwise.

When $v=1$, the group $R'$ is the transposition
of $C_B=\{b,b_1\}$. Suppose a left mark exchanges
these two points. It centralizes the full right
factor. Every right mark after $j+1$ fixes $b$,
so would also fix $b_1$. The first right mark
preserves the pair, as does $L$ by the definition
of $C_B$. Thus the complete $LR$ would preserve
$C_B$, making $B=C_B$. The full right factor
also fixes the other side $A$, as just proved.
It would have no active label after $j+1$,
contrary to $y>x=j+2$. Hence $L$ fixes $C_B$
pointwise in this case as well.
In either case $L$ fixes $b$ and centralizes
$R$, so the other actual side is $B=b^R$, with
$L$ fixing all of $B$. The right action on $B$
is transitive; it need not be symmetric.

Reverse the $h$ left marks on $A$ and use their
exhausted anchored path $a_0=a,\ldots,a_h$.
As in the complete minimum-orbit proof, the
only possible extra transposition in those marks
is $(a_{h-1}\ a_h)$ in the mark three positions
from the end. If present, deleting the first
mark $\rho_m$ would leave the left orbit graph
connected. The crossing and the transitive
right action would then make $Q_{>m}$ transitive
on $W$, a contradiction. Hence the remaining
left marks give exactly $\Sym(A\setminus\{a_h\})$
and fix $a_h$.

For $h\ge3$, $\rho_j$ centralizes the pure
$\Sym(\{a_1,\ldots,a_h\})$, so fixes its support.
For $h=2$, its possible extra transposition
$(a_1\ a_2)$ would again make the first omission
transitive through its connected orbit graph;
therefore it is absent. Thus $z=a_h$ is fixed
by $\rho_j$ and every mark of $Q_{>m}$.
The remaining left marks are transitive on
$A\setminus\{z\}$, the right marks are transitive
on $B$, and the crossing joins these sets.
They give one transitive complementary orbit,
as claimed. In particular $z\in A=C_A$.
\end{proof}

\begin{lemma}[The short left sides]
In the same actual local crossing setting, suppose
$h=j-m$ is zero or one, and the first omission
$Q_{>m}$ is intransitive. Let $R'=Q_{>j}$ and
$v=y-j\ge1$, or let
$R'=Q_{\{j+1,\ldots,x-1\}}$ and
$1\le v=x-j-1<y-j$. Put
$C_A=a^{LR'}$, $C_B=b^{LR'}$ as before.
If $|C_A|+|C_B|=h+v+2$, the first omission has
one fixed point in $C_A\cup C_B$ and one
transitive complementary orbit on $W$.
\end{lemma}
\begin{proof}
The anchored action and strict prefix results
give active right lengths $v_A,v_B$ on these
sets with maximum $v$. An orbit through its
anchor under $R'$ has at least $v_A+1$ or
$v_B+1$ points, respectively.

If $h=0$, the left factor is trivial. Equality
forces one of these endpoint orbits to have
$v+1$ points and the other to be a singleton,
with no active $R'$-mark there. Every remaining
right mark, if there is one, fixes that other
anchor because its index is at least $j+2$.
Thus one actual omission side is a singleton.
The first omission is the full right factor
and is transitive on the other side, proving
the assertion.

Suppose $h=1$, so $L$ is generated by the
single mark $\rho_m$. The full right factor
$R=Q_{>j}$ cannot be transitive on both actual
omission sides: otherwise $\gen{\rho_j,R}=Q_{>m}$
would be transitive on $W$. Name a side where
$R$ is intransitive as $A$. Since $LR$ is
transitive there and $L$ has order two, $R$
has precisely two orbits on $A$, exchanged by
$L$. Consequently $C_A$ is a disjoint union
of two equally sized $R'$-orbits, and
\[
                   |C_A|\ge2(v_A+1),\qquad
                   |C_B|\ge v_B+1.
\]
If $v_A=v$, their total would be at least
$2v+3>v+3$, contradicting equality. Hence
$v_B=v$, and equality forces $v_A=0$,
$|C_A|=2$ and $|C_B|=v+1$.
The group $R'$ fixes both points of $C_A$.
Every further right mark fixes $a$ and commutes
with $L$, so also fixes its partner. Thus
$A=C_A$ and $R$ fixes $A$ pointwise.
On the other side, $b^{R'}$ already has at
least $v+1=|C_B|$ points, so $C_B=b^{R'}$.
The group $L$ therefore sends $b$ into its
$R$-orbit. Commutation now gives $B=b^R$,
with $R$ transitive on $B$.

The crossing joins $a$ to $B$. The other point
$z$ of the two-point side $A$ is fixed by
$\rho_j$: its sole possible partner in $A$
is already used by the crossing, and a second
crossing to $B$ is forbidden by local uniqueness. Therefore
$\gen{\rho_j,R}=Q_{>m}$ fixes $z$ and is
transitive on the complement. This proves
both the full and truncated versions without
assuming that the right action is symmetric.
\end{proof}

Retain now every actual odd-degree head and global
paired-tail hypothesis, and choose the colour $j$
by the complete simultaneous-witness theorem.
Write $H=G_{\ge5}$, $K=G_{>5}$ and
$\sigma=\rho_4$. The first three marks are
supported on $D$, and the acquired actual interface is
\[
 \rho_3=(6\ b)\delta,\qquad H|_{D\cup\{b\}}=1,
 \qquad \sigma|_D=1.
\]
Every $K$-intransitive witness orbit for $j$ has a
genuine local unique crossing and is fixed pointwise
by $\delta$.

\begin{theorem}
Suppose two distinct witness orbits $O_1,O_2$
are $K$-intransitive, with last active labels
$x<y$. If $x\ge j+2$, then
\[
                  x<r-1,\qquad
                  |\operatorname{Fix}(G_{>x})|\ge2x+1.
\]
\end{theorem}
\begin{proof}
The acquired two-witness theorem already gives
$x<r-1$ and at least $2x$ fixed points. Suppose
the desired stronger bound fails, and put
$F=\operatorname{Fix}(G_{>x})$. Then $|F|=2x$.
The complete count exhibited the following disjoint
subsets of $F$: the seven points $D\cup\{b\}$,
the whole orbit $O_1$ of size at least $x-3$,
and the two explicit truncated endpoint orbits
$C_A,C_B$ inside $O_2$, of total size at least
$x-4$. Consequently every bound is an equality,
and, with $C=C_A\cup C_B$,
\[
 F=D\cup\{b\}\ \sqcup\ O_1\ \sqcup\ C,
 \qquad |O_1|=x-3,\qquad |C|=x-4.                 \tag{1}
\]
If $j\ge7$, the complete minimum-orbit rigidity
applies on $O_1$, since both side lengths are
at least two; the truncated equality lemma
applies on $O_2$. If $j=5$ or $6$, the short-left-side
lemma applies in its full version on $O_1$ and
its truncated version on $O_2$. Thus in every
case the $K$-orbits on $O_i$ are one fixed
point $z_i$ and a transitive complement $T_i$,
with $z_2\in C$.
The mark $\rho_y\in K$ is nontrivial on $T_2$
but fixes $T_1$ pointwise.

The acquired actual anchor
\[
                         z_0=b^\sigma
\]
is outside $D\cup\{b\}$ and fixed by $K$.
It belongs to $F$, so (1) forces $z_0=z_1$ or
$z_0=z_2$. The permutation $\sigma$ centralizes
$K$ and hence permutes its actual orbits; it
also preserves $F$, since it centralizes $G_{>x}$.
We derive a contradiction in either case.

If $z_0=z_i$ and $\sigma$ preserves $T_i$,
the proper set $D\cup\{b\}\cup O_i$ is invariant
under every generator of $G$. The head preserves
$D$; $\rho_3$ exchanges $6,b$ and fixes $O_i$
through $\delta$; $\sigma$ exchanges $b,z_i$,
fixes $D$ and preserves $T_i$; and $H$ preserves $O_i$ and
fixes $D,b$. The other witness orbit is outside
this set, so this contradicts transitivity.
Thus the complement associated with $z_0$ must
be moved to a distinct $K$-orbit by $\sigma$.

If $z_0=z_1$, every point of $T_1$ lies in $F$,
so its distinct image orbit lies in $F$ as well.
It cannot meet the $K$-fixed head points or
$z_1,z_2$. By (1) it must therefore meet $T_2$,
and so equal $T_2$. This is impossible: the
commuting bijection $\sigma$ would conjugate
$\rho_y|_{T_1}=1$ to the nonidentity restriction
on $T_2$.

If $z_0=z_2$, choose a point in $C\setminus\{z_2\}$;
this set is nonempty since $|C|=x-4\ge3$.
Its image lies in $F$ and in the distinct orbit
$T_2^\sigma$. It cannot be $K$-fixed and cannot
lie in $O_2$, whose only nontrivial $K$-orbit
is $T_2$. Equation (1) therefore places it in
$T_1$. Hence $T_2^\sigma=T_1$, giving exactly
the same contradiction with $\rho_y$.
Both cases are impossible. Thus $|F|\ge2x+1$.
\end{proof}

\paragraph{Attribution and boundary.}
The complete simultaneous-witness count, minimum
local orbit rigidity and actual head anchor are
acquired. The new steps audit equality in
the truncated count, cover the short left sides,
and use the actual $K$-orbits
and $\rho_4$-action on the entire fixed set.
Every invariant-set assertion checks all generator
families. The cases $x=j$ or $x=j+1$, and a
second witness with transitive $K$-action remain,
as do full source routing and global attachments.
No whole-original, question, numerical or formal
credit is claimed.
\end{document}
